\section{Results and Discussion}\label{sec:results}

\subsection{SCoPE Classification Performance}\label{subsec:perf}

Table~\ref{tab:scope_results} summarises the test performance of SCoPE on the
three benchmarks under the 70/30 test/search protocol. For each dataset, $k$
denotes the support-set size with the highest test AUC-ROC. The table also
reports the \texttt{variable} setting, in which the support size is sampled
uniformly from $\{3,\dots,20\}$ for every query.

\begin{table}[!htb]
\caption{SCoPE classification performance under the 70\% test /
30\% search protocol (random partition, no scaffold stratification).
$k$ is the support-set size of the best test-set configuration
observed per dataset (highest test AUC-ROC); Var.: \texttt{variable} support-set
size ($k\in\{3,\dots,20\}$ sampled per query). AUC-ROC: area under the
ROC curve; F1: macro F1-score; Bal.-Acc.: balanced accuracy; MCC:
Matthews correlation coefficient.}\label{tab:scope_results}
\begin{tabular}{@{}lccccc@{}}
\toprule
Dataset & $k$ & AUC-ROC & F1 & Bal.-Acc. & MCC \\
\midrule
ClinTox & 10   & 0.933 & 0.741 & 0.890 & 0.550 \\
        & Var. & 0.912 & 0.695 & 0.842 & 0.466 \\
\midrule
BACE    &  3   & 0.585 & 0.579 & 0.579 & 0.158 \\
        & Var. & 0.533 & 0.532 & 0.533 & 0.065 \\
\midrule
BBBP    &  8   & 0.778 & 0.711 & 0.739 & 0.438 \\
        & Var. & 0.749 & 0.719 & 0.749 & 0.454 \\
\botrule
\end{tabular}
\end{table}

On ClinTox, SCoPE attains an AUC-ROC of 0.933 at $k=10$ and an accuracy of
88.0\% despite the severe class imbalance ($\approx$8\% positive). The macro
F1 (0.741) and MCC (0.550) are lower than the AUC-ROC and the balanced
accuracy (0.890), which reflects the inherent difficulty of detecting rare
toxic compounds from a small support set, a challenge shared by low-data
approaches \cite{zhang2024molfescue,wu2018moleculenet}. Since AUC-ROC is the headline metric, $k=10$ is
reported in Table~\ref{tab:scope_results}; however, this choice does not
dominate every metric across the ClinTox sweep: the best balanced accuracy
(0.903) is obtained at $k=8$, the best macro F1 (0.756) and MCC (0.570) at
$k=15$, and the AUC-ROC at $k=10$ and $k=15$ is practically tied (0.933).
Furthermore, the size selected on the search split for ClinTox is $k=5$
(CV balanced accuracy 0.946), for which the test AUC-ROC and MCC are 0.884 and
0.498, respectively. Table~\ref{tab:scope_results} should therefore be read as
an upper reference for ClinTox, whereas for BACE ($k=3$) and BBBP ($k=8$) the
best test size coincides with the best search size among the fixed sizes
(Appendix~\ref{app:sweep}).

In all three datasets, the \texttt{variable} setting attains a lower test
AUC-ROC than the best fixed size (0.912 vs.\ 0.933 on ClinTox, 0.533 vs.\ 0.585
on BACE, and 0.749 vs.\ 0.778 on BBBP). On ClinTox its balanced accuracy (0.842)
and MCC (0.466) are also lower than those of the best fixed size. On BBBP,
however, its macro F1 (0.719) is higher than that of the $k=8$ configuration
(0.711) and its MCC (0.454) is within 0.02 of the best MCC of the sweep
(0.474 at $k=15$). These results indicate that a support-set size selected per
dataset is preferable to mixing sizes within a single run.

On BACE, the AUC-ROC of 0.585 (MCC 0.158) is close to chance and lies below the
values reported for the baselines in Section~\ref{subsec:compare}
(Table~\ref{tab:context}), although the protocols differ; this indicates that
the compression signal alone under random 70/30 splitting is insufficient for
this binding-affinity task. On BBBP, the AUC-ROC of 0.778 (MCC 0.438) indicates
that SMILES string compression encodes a partial signal for blood-brain barrier
penetration, a property driven by molecular size, polarity, and hydrogen-bonding
capacity, all of which manifest in SMILES character patterns \cite{martins2012bayesian}.
Across the three best test-set configurations of
Table~\ref{tab:scope_results}, the mean test AUC-ROC amounts to 0.765. These
results support the hypothesis that raw SMILES strings, despite their
non-unique graph encoding \cite{li2022novel,bennett1998information}, contain sufficient structural
redundancy for compression-based dissimilarity to recover useful similarity
rankings, consistent with the theoretical upper bound of NCD on the true NID
\cite{cilibrasi2005clustering}; nevertheless, the gap relative to the baselines underscores the
limits of the raw-string compression signal on structurally heterogeneous
tasks.

The configurations selected by the search at the best test support-set size of
each dataset are listed in Appendix~\ref{app:sweep}. In every dataset and
setting of the sweep, the selected \texttt{keep\_similar} is \texttt{True}. For
the three configurations of Table~\ref{tab:scope_results}, the
compressor--metric subsets vary across tasks: on BACE ($k=3$) the search
retains NCD alone with compressors \{zlib, gzip\}; on BBBP ($k=8$) the full
compressor family \{bz2, zlib, gzip\} is combined with \{CDM, CLM\}; and on
ClinTox ($k=10$), \{bz2, gzip\} with \{CDM, CLM\}. None of these three
configurations relies on a single compressor or activates all three metrics
simultaneously. The two tasks with the strongest test signal (BBBP and ClinTox)
both select CDM and CLM, whereas the weaker BACE signal is captured by NCD
alone. The compressor--metric subsets and \texttt{keep\_similar} are the
\texttt{arg\,max} of the search objective (balanced accuracy, 30\% search
split) over the 98 configurations enumerated per support-set size; only the
choice of $k$ reported in Table~\ref{tab:scope_results} is made on the test
split.

% [FIGURE PLACEHOLDER: classification results (BACE, BBBP, ClinTox) -- add here
% and cite it in the text once the images are available]
